%% 
%% Formatted for Elsevier CAS-SC
%% 

\documentclass[a4paper,fleqn]{cas-dc}

\usepackage[authoryear,longnamesfirst]{natbib}
\usepackage{amsmath,amssymb,amsfonts}
\usepackage{algorithmic}
\usepackage{graphicx}
\usepackage{textcomp}
\usepackage{colortbl}
\usepackage{algorithm,algorithmic}
\usepackage{hyperref}
\hypersetup{hidelinks=true}
\usepackage{multirow, array, makecell}
\usepackage{verbatim}
\usepackage[flushleft]{threeparttable}
\usepackage{booktabs}
\usepackage{tabularx}
\usepackage{adjustbox}

%%%Author macros
\def\tsc#1{\csdef{#1}{\textsc{\lowercase{#1}}\xspace}}
\tsc{WGM}
\tsc{QE}
%%%

\begin{document}
\let\WriteBookmarks\relax
\def\floatpagepagefraction{1}
\def\textpagefraction{.001}

% Short title
\shorttitle{Classification of Motor Imagery EEG Signals}    

% Short author
\shortauthors{Cuascota et al.}  

% Main title of the paper
\title [mode = title]{Classification of Motor Imagery EEG Signals Using CNN-Transformer Hybrid Models}  

% Title footnote 1.
%\tnotemark[1]
%\tnotetext[1]{This work was submitted for review on March 3, 2026. The Universidad San Francisco de Quito USFQ supported this work through the Poli-Grants Program under Grants 33609 and 41993.}

\author[1]{Joel Cuascota}[orcid=0009-0002-4316-5756]
\credit{Software, Validation, Formal analysis, Investigation, Data curation, Writing - original draft, Visualization}

\author[1]{Noel P\'erez-P\'erez}[orcid=0000-0003-3166-745X]
\cormark[1]
\credit{Conceptualization, Methodology, Validation, Formal Analysis, Investigation, Supervision, Writing - review \& editing, Project administration, Funding acquisition}

\author[2]{Jorge Oliveira}[orcid=0000-0002-3190-8367]
\credit{Formal analysis, Writing - review \& editing}

\author[3]{Maria Baldeon-Calisto}[orcid=0000-0001-9379-8151]
\credit{Formal analysis, Writing - review \& editing}

\author[4]{Nathaly Orozco Garzón}[orcid=0000-0002-5232-7529]
\credit{Formal analysis, Writing - review \& editing}

\author[1]{Henry Carvajal Mora}[orcid=0000-0003-0529-8224]
\credit{Formal analysis, Writing - review \& editing}

\author[1]{José Vega-Sánchez}[orcid=0000-0002-3305-1109]
\credit{Formal analysis, Writing - review \& editing}

\author[5]{Edgar A. Marquez}[orcid=0000-0002-7503-1528]
\credit{Visualization, Writing - review \& editing}

% Address/affiliation
\affiliation[1]{organization={Colegio de Ciencias e Ingenier\'ias, Universidad San Francisco de Quito (USFQ)},
            city={Quito},
            postcode={170157},
            country={Ecuador}}

\affiliation[2]{organization={Universidade da Maia and LIACC – Laboratory of Artificial Intelligence and Computer Science},
            city={Maia},
            postcode={4475-690},
            country={Portugal}}

\affiliation[3]{organization={School of Business, Wake Forest University},
            city={Winston-Salem},
            postcode={NC 27109},
            country={USA}}

\affiliation[4]{organization={ETEL Research Group, Faculty of Engineering and Applied Sciences, Networking and Telecommunications Engineering, Universidad de Las Américas (UDLA)},
            city={Quito},
            postcode={170503},
            country={Ecuador}}

\affiliation[5]{organization={Grupo de Investigaciones en Química y Biología, Departamento de Química y Biología, Facultad de Ciencias Exactas, Universidad del Norte},
            addressline={Carrera 51B, Km 5, vía Puerto Colombia},
            city={Barranquilla},
            postcode={081007},
            country={Colombia}}

% Corresponding author text
\cortext[1]{Corresponding author: nperez@usfq.edu.ec}

% Here goes the abstract
\begin{abstract}
Motor imagery-based brain-computer interfaces aim to translate electroencephalographic activity into control commands without requiring overt movement. However, strong inter-subject variability, exacerbated by the limited spatial resolution of compact electrode montages, makes subject-independent decoding particularly challenging. This work presents a compact hybrid convolutional neural network–Transformer architecture designed for stable binary left/right motor imagery decoding using a low-density eight-channel sensorimotor montage. The model integrates convolutional feature extraction with a parameter-efficient temporal self-attention mechanism to capture both localized spatiotemporal patterns and long-range temporal dependencies in the electroencephalographic signal. Under strict subject-wise five-fold cross-validation on the PhysioNet EEG Motor Movement/Imagery dataset, the proposed convolutional neural network-Transformer achieved a mean macro F1-score of 0.835. Moreover, the statistical analysis revealed no significant performance difference ($p > 0.05$) relative to larger architectural variants, thus justifying the selection of the proposed compact model for its computational efficiency and scalability. With only 0.232 million parameters and approximately 0.024 giga-floating-point operations per forward pass, the proposed convolutional neural network-Transformer constitutes a compact and resource-efficient backbone that supports low-latency subject-independent motor imagery decoding under low-density EEG.
\end{abstract}

% Research highlights
\begin{highlights}
\item Compact CNN–Transformer decodes MI-EEG using eight channels.
\item Subject-wise validation achieved 0.835 macro F1-score.
\item Attention and saliency maps support physiological plausibility.
\item The model uses 0.232 M parameters and 0.024 GFLOPs.
\item Low-density EEG decoding supports wearable BCI development.
\end{highlights}

% Keywords
\begin{keywords}
Brain-computer interfaces \sep Electroencephalography \sep Motor imagery \sep Cross-subject decoding \sep Convolutional neural networks
\end{keywords}

\maketitle

%==============================================================
\section{Introduction}
\label{sec:introduction}
%==============================================================

Brain-computer interfaces (BCIs) seek to translate brain activity into digital commands to enable communication and control without reliance on peripheral nerves or muscles \citep{1niedermeyer2005electroencephalography}. Among the most studied paradigms is motor imagery (MI), in which users imagine limb movements such as left-hand (L) or right-hand (R) motions. MI-BCIs hold promise for neurorehabilitation, prosthetic control, and assistive communication for individuals with severe motor impairments \citep{2pfurtscheller1999event,3schalk2004bci2000}. EEG signals recorded during MI exhibit characteristic wave frequency patterns in the $\mu$ (8-12~Hz) and $\beta$ (13-30~Hz) bands over the primary motor cortex (C3, Cz, C4) \citep{4pfurtscheller2001functional}. However, decoding MI from scalp EEG remains challenging. Due to a low signal-to-noise ratio, physiological and environmental artifacts are commonly observed in the EEG trace, compounding the pronounced intra- and inter-subject variability \citep{5physionet2009eeg,6dhiman2023ml_eeg_bci, SahaBaumert2020_IntraInterSubjectEEG}.
 
Despite decades of progress, EEG-based BCIs continue to exhibit substantial performance variability across participants, underscoring a persistent challenge for practical MI decoding. Differences in cortical activation patterns, electrode placement, and noise characteristics often weaken model generalization, a limitation that becomes even more pronounced under low-density electrode montages, where reduced spatial coverage limits class discriminability \citep{Xu2020CrossDataset, abdullah2022eeg}. This difficulty is further compounded by the relatively small number of trials available per subject in many MI datasets, which restricts the effective training of deep neural networks \citep{Lotte2018EEGReview}. As a result, an important open problem remains the design of MI decoding frameworks that can jointly maintain competitive performance under compact electrode configurations, generalize across unseen subjects, and remain computationally lightweight enough to support translational deployment scenarios. Existing approaches often improve one or two of these dimensions, but typically at the expense of the others.

For example, existing MI-EEG decoders face important limitations. Many deep models rely on dense montages (32-64 channels) to achieve accurate results, but under compact configurations such as an eight-channel motor-area setup, a drop in their accuracy has been observed \citep{nagarajan2023relevance}. Also, cross-subject generalization remains particularly challenging, with performance dropping substantially when models are evaluated on unseen individuals. Moreover, several architectures incur high computational costs, limiting their applicability in real-time or wearable BCI systems \citep{wang2020accurate}. This unresolved trade-off strongly motivates the development of unified architectures that can jointly maintain stability under low-density montages (spatial sparsity), generalize well across subjects (inter-subject variability), and maintain low computational complexity (efficiency constraints), thereby advancing MI-based BCI systems toward reliable real-world deployment.


Therefore, we propose a compact and computationally efficient CNN–Transformer hybrid architecture for subject-independent decoding of left/right motor imagery on the PhysioNet EEG dataset, specifically designed to operate with a low-density 8-channel sensorimotor montage while addressing the joint challenges of spatial sparsity, inter-subject variability, and deployment-oriented efficiency. Rather than introducing a Transformer-based architecture in isolation, the present work focuses on demonstrating that a carefully constrained hybrid design can offer a favorable balance among cross-subject decoding performance, compactness, and calibration-free operation in a realistic low-density setting. The main contributions behind the goal are related to:

\begin{itemize}
    \item A compact hybrid CNN–Transformer architecture for low-density MI-EEG decoding: we introduce a lightweight CNN–Transformer backbone that couples depthwise-separable convolutional blocks with a temporal self-attention encoder to jointly model local sensorimotor patterns and long-range temporal dependencies in 8-channel MI-EEG signals. This design is explicitly tailored to subject-independent L/R decoding under spatially sparse electrode configurations.
    \item A subject-independent training and model-selection framework for low-density EEG: we formulate and evaluate the proposed decoder under a strict subject-wise five-fold cross-validation protocol, combined with controlled architectural exploration and statistically grounded model selection. This framework is designed to identify architectures that remain stable across unseen subjects while avoiding unnecessary model complexity.
    
    \item An efficiency-oriented MI decoding pipeline suitable for translational BCI settings: we develop an end-to-end low-density decoding pipeline that integrates minimal preprocessing, lightweight augmentation, compact feature extraction, and a single-layer Transformer encoder, yielding a resource-constrained design appropriate for real-time and wearable MI-BCI scenarios.
\end{itemize}

The rest of this work is organized as follows. Subsection~\ref{subsec:related_work} reviews the related literature. Section~II presents the Materials and Methods, including the PhysioNet dataset and the proposed architecture. Section~III reports the experimental results, and Section~IV summarizes the conclusions.


%==============================================================
\subsection{Related work}
\label{subsec:related_work}
%==============================================================

Over the past decade, numerous approaches have been reported in the literature for classifying MI episodes, evolving from traditional machine learning to deep neural networks as presented next:

\subsubsection{Traditional machine learning approaches}
Early MI decoders primarily relied on spatial filtering and shallow learning techniques. For example, Ang et al. introduced the Filter Bank Common Spatial Pattern (FBCSP) algorithm, which projects EEGs into discriminative spatial components using multiple frequency bands \citep{7ang2008filter}. This method was validated on the BCI Competition IV datasets (2a and 2b), achieving cross-validation accuracies of approximately 90\% in subject-dependent settings. More recently, Xiong et al. proposed an improved FBCSP variant that employs Particle Swarm Optimization (PSO) to adaptively select spatial filters \citep{14xiong2023improved}. When evaluated on the BCI Competition IV 2a and 2b datasets, this method reached classification accuracies of 74.61\% and 81.19\%, respectively, for four-class MI. Additionally, Barachant et al. applied Riemannian geometry to classify covariance matrices directly on the tangent space \citep{8barachant2011multiclass}. Their approach, validated on the BCI Competition IV 2a dataset, outperformed standard pipelines based on Common Spatial Pattern (CSP) and Linear Discriminant Analysis (LDA), increasing the mean classification accuracy from 65.1\% to 70.2\%.

\subsubsection{Deep learning and CNNs}
Deep learning approaches have achieved competitive performance by automating feature extraction. For example, Lawhern et al. introduced EEGNet, a compact convolutional neural network designed for general EEG decoding \citep{9lawhern2018eegnet}. While highly efficient, its performance varies with montage density. Building on convolutional architectures, Ding et al. proposed TSception, which extends multi-scale temporal convolutions to capture diverse frequency bands \citep{10ding2020tsceptiondeeplearningframework}. They reported an accuracy of 86.03\% on an 18-subject VR-based emotion EEG dataset. Similarly, Zhang et al. developed EEG-Inception, a deep CNN featuring parallel inception modules \citep{zhang2021eeg}. On the BCI Competition IV 2b dataset, it achieved a global accuracy of 77.5\%, which improved to 88.6\% after subject-specific fine-tuning. In a multi-class setting, Rao et al. optimized an EEGNet variant for decoding four-class finger flexion, reporting a classification accuracy of 72.91\% on a proprietary dataset \citep{15rao2024optimized}.

To mitigate data scarcity, advanced training strategies have been explored. Banville et al. demonstrated that self-supervised pretraining on large unlabeled sleep EEG datasets led to consistent performance improvements in downstream pathology detection tasks \citep{12banville2019selfsupervisedrepresentationlearningelectroencephalography}. In the context of emotion recognition, Hu et al. applied contrastive self-supervised learning to the Shanghai Jiao Tong University (SJTU) Emotion EEG Dataset (SEED) and its four-emotion extension (SEED-IV), achieving significant accuracy gains under limited label availability \citep{11hu2024contrastive}. For motor imagery, Han et al. proposed a semi-supervised contrastive learning framework evaluated on the BCI Competition IV 2a dataset \citep{13han2021semi}. Their method achieved a subject-independent accuracy of 50.8\%, outperforming several supervised CNN baselines in low-data regimes.

\subsubsection{Transformers and hybrid architectures}
Recently, Transformer-based architectures have been adopted to model long-range global dependencies that CNNs may miss. For example, Zhang et al. proposed a local-global convolutional transformer that captures both temporal dynamics and spatial discrepancies between hemispheres \citep{zhang2023local}. Their method achieved improved decoding accuracy on a publicly available Korean MI dataset compared with standard CNNs. Furthermore, Zhao et al. introduced the Improved Multi-scale Convolution and Transformer Network (IMCTNet), which integrates channel attention with a transformer encoder \citep{zhao2025multi}. This architecture achieved superior classification accuracies of 81.83\% on the BCI Competition IV 2a dataset and 86.47\% on dataset 2b. Similarly, Wu et al. developed AMEEGNet, an attention-based multiscale EEGNet \citep{wu2025ameegnet}. This model reported state-of-the-art accuracies of 81.17\% on BCI-2a and 89.83\% on BCI-2b, demonstrating the efficacy of combining attention mechanisms with multi-scale feature extraction. 

Despite the significant progress achieved by traditional, deep learning, and hybrid Transformer-based approaches, several key limitations remain. Many published methods are evaluated in subject-dependent settings, rely on dense montages, or require subject-specific calibration, fine-tuning, or online adaptation to achieve their best performance. Furthermore, while recent hybrid architectures have shown the value of attention mechanisms for EEG decoding, these models are often not explicitly optimized for sparse-channel operation or for deployment under tight computational constraints. In this context, the contribution of the present work is not merely the use of a CNN–Transformer combination, but the demonstration that a compact hybrid architecture can sustain competitive subject-independent performance using only eight channels, without test-time adaptation, while maintaining a low computational footprint. This deployment-oriented framing differentiates the proposed model from prior work that prioritizes accuracy without simultaneously addressing montage sparsity, cross-subject generalization, and efficiency.


%==============================================================
\section{Materials and methods}
%==============================================================

\subsection{Database}
We used the public EEG Motor Movement/Imagery (EEGMMI) dataset from PhysioNet \citep{5physionet2009eeg}, recorded with the BCI2000 acquisition system \citep{3schalk2004bci2000}. The database contains EEG recordings from 109 healthy volunteers acquired with a 64-channel 10-10 montage at 160 Hz. Each participant completed 14 runs: 2 resting-state baselines and 12 motor-related runs, including both motor execution (ME) and MI tasks involving unilateral and bilateral hand and foot movements. Because the runs are organized as continuous blocks rather than discrete trials, EEGMMI is well-suited for window-based MI decoding and for evaluating subject-independent performance under realistic inter-subject variability. A random example of L- and R-hand motor imagery epochs over C3, C4, and Cz is shown in Fig.~\ref{fig:raw_signals}, illustrating the subtle lateralized $\mu/\beta$ modulations exploited by the proposed model. 

\begin{figure*}[!t]
    \centering
    \includegraphics[width=0.65\textwidth]{fig1_raw_signals.png}
    \caption{Single-trial EEG example of Left (blue) and Right (red) motor imagery from a randomly selected subject, displayed on three sensorimotor channels (C3, C4, Cz). The dashed vertical line indicates cue onset ($t=0$).}
    \label{fig:raw_signals}
   % \vspace{-10pt}
\end{figure*}

\subsection{Proposed method}
The proposed CNN-Transformer is a hybrid architecture that combines a convolutional front-end with a compact Transformer encoder to model both local and long-range temporal structure in MI-EEG. This design leverages the strengths of each component: convolutions extract consistent sensorimotor rhythms and short-range temporal patterns, while self-attention enables the model to integrate information across the entire signal window, as shown in Fig.~\ref{fig:cnn_transformer}. 

From this figure, the input consists of multichannel EEG segments from the eight motor-area electrodes, represented as a matrix $X \in \mathbb{R}^{C \times T}$, where $C$ is the number of channels ($C=8$) and $T$ is the number of time samples ($T=960$) (see Fig.~\ref{fig:cnn_transformer}, step 1). A sequence of convolutional blocks first transforms the raw signal into higher-level spatiotemporal embeddings. The stem layer performs temporal filtering with a wide kernel ($k=129$, stride $s=2$) to capture low-frequency $\mu$- and $\beta$-band modulations. Subsequently, two depthwise-separable convolutional blocks enable efficient channel-wise filtering and cross-channel feature mixing with substantially fewer parameters. Each block is followed by GroupNorm, ELU activations, and dropout with a rate of $0.2$. The GroupNorm normalizes features independently of the batch size, providing stable optimization and improved generalization for EEG data, which typically exhibits high inter-subject variance. A final $1\times1$ convolution projects the feature maps into a fixed-dimensional latent sequence of length $L=144$ (see Fig.~\ref{fig:cnn_transformer}, step 2).

This $L$-sequence is then processed by a single-layer Transformer encoder, which previously adds positional encodings and a learnable classification token (\texttt{[CLS]}) to preserve temporal order and aggregate global information in a dedicated vector, respectively. 
The multi-head self-attention module computes data-dependent attention scores for temporal interactions \citep{vaswani2017attention}, allowing the model to flexibly aggregate features across time and integrate short transients with sustained ERD/ERS patterns. Importantly, this adaptive weighting mechanism enables the model to emphasize the most informative temporal segments while suppressing irrelevant or noisy components, which is particularly beneficial under low-density electrode configurations. In addition, by learning subject-invariant temporal dependencies, the attention mechanism improves stability against inter-subject variability, a key challenge in practical MI-EEG decoding. Here, ERD and event-related synchronization (ERS) refer to the task-induced decreases or increases, respectively, in the power of sensorimotor rhythms, particularly in the $\mu$ and $\beta$ bands, during motor imagery. A feed-forward network refines the attended representation, with residual connections and layer normalization in both sublayers. After encoding, the final \texttt{[CLS]} vector is normalized and passed through a linear layer to produce the L/R class prediction (see Fig.~\ref{fig:cnn_transformer}, steps~2~and~3).

During model development, several architectural variants were explored by varying the depth of the convolutional encoder and the number of attention heads, while keeping the training protocol and preprocessing fixed, as summarized in Table~\ref{tab:arch_space}. These controlled variations enabled a systematic evaluation of 9 architectures to analyze how convolutional capacity and attention complexity influence subject-independent performance and to identify the best-performing architecture.

% =============================
% TABLE: ARCHITECTURE SWEEP
% =============================
\begin{table}[h]
\centering
\caption{Architectural variants explored during the model design stage.}
\label{tab:arch_space}
\footnotesize
\begin{threeparttable}
\resizebox{\columnwidth}{!}{
\begin{tabular}{ccc}
\toprule
$n_{\text{dw}}$ & Attention Heads & Notes \\
\midrule
0 & \{2, 4, 6\} & Shallow CNN front-end \\
2 & \{2, 4, 6\} & Moderate convolutional capacity \\
4 & \{2, 4, 6\} & Deepest convolutional front-end \\
\bottomrule
\end{tabular}
}
\begin{tablenotes}
\begin{minipage}{\columnwidth}
    \item  $n_{\text{dw}}$ denotes the number of depthwise-separable convolutional blocks in the convolutional front-end.
\end{minipage}
\end{tablenotes}
\end{threeparttable}
%\vspace{-10pt}
\end{table}

\begin{figure*}[t!]
    \centering
    \includegraphics[width=\textwidth]{fig2_EEG_workflow.png}
    \caption{Architecture of the proposed hybrid CNN-Transformer model for EEG-based left/right hand classification. First, an 8-channel EEG segment from a 6s window is processed by a temporal Conv1D followed by depthwise separable convolution blocks with GroupNorm, ELU activation, and dropout. The resulting feature maps are projected with a 1×1 Conv1D and converted into a sequence with positional encoding and a prepended [CLS] token. Second, a Transformer encoder models temporal dependencies, and the final [CLS] representation is passed through LayerNorm and a linear classifier to predict left/right hand class labels.}
    \label{fig:cnn_transformer}
    %\vspace{-10pt}
\end{figure*}

\subsection{Experimental setup}

\subsubsection{Experimental dataset creation}
Although the PhysioNet EEG Motor Movement/Imagery database includes 109 volunteers, several recordings contain incomplete sessions, missing labels, or corrupted EEG signals along the 64-channel montage. Thus, we extracted a curated subset of 103 subjects for our analysis, following the quality-control criteria used in \citep{fukushima2024spatiotemporal}. From this subset, we created a balanced binary dataset named \textit{PhysioNet-MI-2C}, containing 4,326 valid motor imagery trials, distributed as 2,163 L and 2,163 R trials. 
The binary L/R decoding is well-characterized physiologically and widely used in MI-BCI systems due to the distinct contralateral modulation of sensorimotor rhythms. This curated and balanced dataset design ensures a controlled and reliable evaluation framework, minimizing biases introduced by noisy or incomplete recordings. Moreover, by enforcing strict trial selection and class balance, the resulting dataset provides a consistent benchmark for assessing model performance under subject-independent conditions, which is critical for fair comparison and reproducibility in MI-EEG studies.

Additionally, the signals were reduced from 64 channels to 8 channels using the montage \{C3, C4, Cz, FC3, FC4, CP3, CP4, FCz\}, targeting the core sensorimotor strip and adjacent premotor and somatosensory regions. These sites capture the dominant $\mu$ (8-12~Hz) and $\beta$ (13-30~Hz) ERD/ERS dynamics associated with motor imagery \citep{tsuchimoto2021use}, making them an anatomically grounded choice for low-density MI decoding.

\subsubsection{Signal processing}
All signals were preprocessed using a 60~Hz notch filter to remove line noise, following evidence that convolutional encoders can learn frequency-selective features and attenuate common artifacts without heavy preprocessing \citep{9lawhern2018eegnet,widmann2015digital}. Subsequently, motor imagery segments were extracted using a 6s window from $-1$ to $5$~s relative to cue onset, capturing preparatory and sustained ERD/ERS dynamics characteristic of MI \citep{tibor2017deep}, as shown in Fig.~\ref{fig:windowing}. Finally, each EEG channel was standardized independently using its respective channel-wise mean and standard deviation computed from the training data, and the same normalization parameters were applied consistently to all validation and test signals. As illustrated in Fig.~\ref{fig:windowing}, the epoch extraction process analyzes the continuous multi-channel EEG recording, which is segmented around each motor imagery cue, yielding fixed-size input tensors of dimension $8 \times 960$ (8~channels $\times$ 960~time samples at 160~Hz). 

\begin{figure*}
    \centering
    \includegraphics[width=0.8\textwidth]{fig3_data_extraction.png}
    \caption{Epoch extraction from continuous EEG recordings. A 6s analysis window, spanning from $-1$s to $+5$s relative to the cue onset, is selected from the continuous 8-channel EEG signal. The highlighted segment is then extracted as a fixed-length epoch containing the channels C3, C4, Cz, CP3, CP4, FC3, FC4, and FCz. Each resulting epoch forms one input sample for the proposed model, represented as an $8\times960$ tensor.}
    \label{fig:windowing}
    %\vspace{-10pt}
\end{figure*}

\subsubsection{Training, validation, and test Sets}

All experiments followed a strict subject-wise five-fold cross-validation scheme to prevent cross-participant leakage. The full set of 103 subjects from \textit{PhysioNet-MI-2C} was used in a rotational manner: in each fold, 20–21 subjects were reserved as the hold-out test set, while the remaining 82–83 subjects formed the development pool. Within this pool, subjects were further split at the participant level into training (67–68 subjects) and validation (15 subjects) sets, ensuring that all samples from a given subject remained in a single partition per fold, as shown in Fig.~\ref{fig:kfold_split}. 

Importantly, for each of the five folds, a completely independent model instance was initialized with random weights and trained from scratch using exclusively its designated training and validation pools, ensuring no parameter transfer or information leakage across folds.

Additionally, to improve cross-subject robustness, data augmentation was applied online to the training set only, including circular time shifts up to 3\% of the window length ($\pm0.18 s$), additive white Gaussian noise following $N(0, 0.03^2)$, and random masking of at most one of the eight channels. These perturbations preserve the ERD/ERS structure while following established EEG augmentation practices reported to improve generalization without distorting class-relevant information \citep{zhang2020data,george2022data,he2021data}. 

\begin{figure}
    \centering
    \includegraphics[width=\columnwidth]{fig4_patients_distribution.png}
    \caption{Subject-wise five-fold cross-validation scheme used in all experiments. Subjects are strictly separated across the training, validation, and hold-out test sets.}
    \label{fig:kfold_split}
   % \vspace{-10pt}
\end{figure}

% --------------------------------------------------
\subsubsection{Model optimization}
\label{model_optimization}
Once the best-performing architecture was selected as the fixed backbone, the proposed CNN–Transformer was optimized using a restricted grid search over key training hyperparameters (see Table~\ref{tab:hyperparams_tuning}). The search space included the base learning rate, dropout rate in the convolutional and Transformer blocks, and the number of Transformer encoder layers, while the embedding dimension ($d=144$), warmup epochs ($4$), and early-stopping patience ($12$) were kept fixed. Each configuration was evaluated under the full subject-wise five-fold cross-validation protocol. During training, \textit{AdamW} with a warmup–cosine learning-rate schedule was used for optimization, focal loss with inverse-frequency weighting was adopted to handle class imbalance, a balanced per-subject sampler was employed to reduce inter-subject bias, and an exponential moving average (EMA) of the model parameters was used to improve validation and test stability.

% =============================
% TABLE: HYPERPARAMETER SEARCH
% =============================
\begin{table}[t]
\centering
\caption{Hyperparameter search space explored during the statistical fine-tuning phase.}
\label{tab:hyperparams_tuning}
\footnotesize
\begin{threeparttable}
\resizebox{\columnwidth}{!}{
\begin{tabular}{l c c}
\toprule
Hyperparameter & Symbol & Search Values \\
\midrule
Encoder layers & $L$ & \{1, 2\} \\
Base learning rate & $\eta_{\text{base}}$ & \{$5\times10^{-4}$, $1\times10^{-3}$\} \\
Dropout rate (CNN \& Enc) & $p_{\text{drop}}$ & \{0.1, 0.2\} \\
\bottomrule
\end{tabular}
}
\end{threeparttable}
%\vspace{-10pt}
\end{table}

\subsubsection{Baseline methods}
To ensure a rigorous and controlled comparison against established architectures, two baseline models were evaluated under conditions strictly identical to those used for the proposed model. The first baseline is EEGNet \citep{lawhern2018eegnet}, a highly efficient and widely used convolutional architecture designed specifically for EEG decoding. The standard EEGNet configuration was implemented, utilizing 16 temporal filters and 4 spatial depthwise filters, yielding a compact model with 7,010 parameters. The second baseline is a CNN-only variant created by ablating the Transformer encoder from the proposed architecture. This model preserves the identical convolutional front-end (a 129-kernel temporal stem and two depthwise-separable blocks) but replaces the attention mechanism and the \texttt{[CLS]} token with a global average pooling layer, resulting in 46,178 parameters. Both baselines were trained and tested using the exact same subject-wise five-fold splits, channels, preprocessing, and optimization strategy to guarantee a fair comparison focused strictly on architectural capacity.

\subsubsection{Assessment metrics and selection rule}
\label{model_selection_rule}
The overall performance of the proposed architecture was obtained by evaluating each of the five independent models solely on its respective unseen hold-out test set, and subsequently averaging these results. The metrics computed across the five folds included accuracy (ACC), precision (PRE), recall (REC), macro F1-score, specificity (SPEC), and sensitivity (SENS). Among these metrics, the mean macro F1-score was selected as the primary criterion because it provides a balanced measure of class separability in the binary L/R task. For model selection, the configurations were first ranked according to mean macro F1-score across folds. The fold-wise Wilcoxon signed-rank tests were then used to compare the top-performing configuration against more efficient alternatives. Because the number of folds is limited, these tests were used as support for model selection rather than as definitive evidence of universal superiority. Subsequently, following a parsimony criterion, when no statistically significant difference ($p>0.05$) was found, the lower-complexity model was preferred. Besides, footprint and runtime measures, including parameter count, inference latency, and FLOPs, were also reported to assess scalability.

\subsubsection{Platform Implementation}
All methods were implemented in Python 3.11 and executed within a Linux container on a high-performance workstation equipped with NVIDIA H200 GPUs and Intel Xeon Platinum processors. Experiments were run on a single H200 GPU with approximately 144 GB of memory. Deep-learning models were implemented in PyTorch, and EEG preprocessing was performed with MNE-Python. To ensure reproducibility, preprocessing, training, and cross-validation experiments were conducted using a fixed random seed of 42. 


%==============================================================
\section{Results and discussion}
\label{sec:results}
%==============================================================
This section evaluates the proposed CNN–Transformer for subject-independent MI decoding under the adopted subject-wise five-fold cross-validation protocol. Quantitative results are presented first, followed by qualitative analyses of saliency and self-attention patterns to interpret the learned representations. For clarity, models are denoted as nb$X$\_h$Y$, where $X$ is the number of depthwise convolutional blocks in the CNN front-end, and $Y$ is the number of attention heads in the Transformer encoder.

\subsection{Cross-validation architecture evaluation}
\label{subsec:train_performance}

The nine CNN-Transformer variants evaluated in the binary L/R task achieved consistently good performance (mean macro F1-score values between 0.808 and 0.823). This indicates that even shallow convolutional stems, when paired with compact self-attention, can extract discriminative sensorimotor patterns from low-density eight-channel EEG, as shown in Table~\ref{tab:arch_sweep}.

Within this family of models, the \textit{nb2\_h6} configuration achieved the highest mean macro F1-score of $0.822 \pm 0.016$, combining a moderate convolutional depth ($n_{\text{dw}}=2$) with six attention heads. Nevertheless, deeper configurations, such as \textit{nb4\_h6}, exhibited comparable or slightly lower behavior ($0.809 \pm 0.021$). Statistical assessment using the Wilcoxon signed-rank test confirmed that increasing the model depth offers no significant improvement ($p > 0.05$). This trend is consistent with expectations for low-density MI-EEG, where the limited spatial sampling restricts the benefits of deeper convolutional hierarchies. Although performance was statistically similar across the top variants, \textit{nb2\_h6} achieved the lowest standard deviation, indicating superior stability across subjects. Furthermore, it maintains a compact parameter footprint (0.232~M) and a low computational cost (0.024~GFLOPs). Thus, according to the selection criteria, it is considered the most appropriate backbone for the task under analysis (see Table~\ref{tab:arch_sweep}, bold line).

% =============================
% TABLE: ARCHITECTURE SWEEP
% =============================
\begin{table*}[t]
\centering
\caption{Cross-subject performance of the CNN-Transformer variants under subject-wise five-fold cross-validation.}
\label{tab:arch_sweep}
\footnotesize
\begin{threeparttable}
\resizebox{\textwidth}{!}{
\begin{tabular}{l c c c c c c c c c}
\toprule
Model & Params (M) & Epoch (SD) & ACC (SD) & PRE (SD) & REC (SD) & F1 (SD) & SPEC (SD) & SENS (SD) & $p$-value$^{\dagger}$ \\
\midrule
% --- Bloque nb0 ---
nb0\_h2 & 0.206 & 23.0 (3.6) & 0.810 (0.032) & 0.811 (0.032) & 0.810 (0.032) & 0.810 (0.032) & 0.810 (0.032) & 0.810 (0.032) & 0.0770 \\
nb0\_h4 & 0.206 & 26.6 (2.9) & 0.822 (0.019) & 0.822 (0.019) & 0.822 (0.019) & 0.822 (0.019) & 0.822 (0.019) & 0.822 (0.019) & 0.9786 \\
nb0\_h6 & 0.206 & 21.0 (4.5) & 0.813 (0.020) & 0.816 (0.020) & 0.813 (0.020) & 0.813 (0.020) & 0.813 (0.020) & 0.813 (0.020) & 0.2815 \\
\midrule
% --- Bloque nb2 ---
nb2\_h2 & 0.232 & 25.6 (4.9) & 0.808 (0.023) & 0.809 (0.023) & 0.808 (0.023) & 0.808 (0.023) & 0.808 (0.023) & 0.808 (0.023) & 0.2244 \\
nb2\_h4 & 0.232 & 21.6 (5.5) & 0.808 (0.023) & 0.808 (0.023) & 0.808 (0.023) & 0.808 (0.023) & 0.808 (0.023) & 0.808 (0.023) & 0.0822 \\
\rowcolor[gray]{0.9}
\textbf{nb2\_h6} & \textbf{0.232} & \textbf{26.2 (10.0)} & \textbf{0.823 (0.015)} & \textbf{0.824 (0.015)} & \textbf{0.823 (0.015)} & \textbf{0.822 (0.016)} & \textbf{0.823 (0.015)} & \textbf{0.823 (0.015)} & -- \\
\midrule
% --- Bloque nb4 ---
nb4\_h2 & 0.353 & 30.0 (4.2) & 0.808 (0.020) & 0.808 (0.020) & 0.808 (0.020) & 0.808 (0.020) & 0.808 (0.020) & 0.808 (0.020) & 0.2068 \\
nb4\_h4 & 0.353 & 24.4 (3.5) & 0.808 (0.022) & 0.808 (0.022) & 0.808 (0.022) & 0.808 (0.022) & 0.808 (0.022) & 0.808 (0.022) & 0.2040 \\
nb4\_h6 & 0.353 & 21.8 (3.5) & 0.809 (0.023) & 0.810 (0.023) & 0.809 (0.023) & 0.809 (0.023) & 0.809 (0.023) & 0.809 (0.023) & 0.1976 \\
\bottomrule
\end{tabular}
}
\begin{tablenotes}
\begin{minipage}{\textwidth}
    \item $^{\dagger}$ Two-sided Wilcoxon signed-rank test against \textit{nb2\_h6} ($N=5, \alpha=0.05$). Fixed hyperparameters: $d=144$, $L=1$, LR=$10^{-3}$, Dropout=0.2, Warmup=4. The bold line denotes the best-selected model.
\end{minipage}
\end{tablenotes}
\end{threeparttable}
%\vspace{-10pt}
\end{table*}

We further optimized the selected \textit{nb2\_h6} backbone by conducting a statistical grid search to refine the Transformer encoder and training hyperparameters. To ensure a comprehensive assessment of stability, each configuration was evaluated using the full five-fold cross-validation protocol. The results are summarized in Table~\ref{tab:grid_search_results}. From this table, it is possible to notice that the best performance was obtained by \textit{Config-07}, which achieved a mean macro F1-score of $0.835 \pm 0.015$. This configuration applies an encoder layer with $L=1$, a learning rate of $1\times10^{-3}$, and a dropout rate of $0.2$. Notably, \textit{Config-07} outperformed deeper variants such as \textit{Config-08} with $L=2$, which reached a lower mean macro F1-score of $0.831$. 

Regarding the increase in the number of encoder layers, statistical testing confirmed that adding a second encoder layer provides no significant advantage when comparing \textit{Config-07} with \textit{Config-08} ($p\text{-value}=0.46$), while substantially increasing the parameter count from 232k to 400k. Furthermore, \textit{Config-07} demonstrated the lowest standard deviation among all tested variations, reinforcing the finding that a compact, single-layer architecture offers the best trade-off between generalization capability and model complexity. Therefore, the \textit{Config-07} configuration was adopted as the final model for all subsequent analyses.

Additionally, we examined the model's learning behavior during optimization, and the average training and validation loss curves across all folds are shown in Fig.~\ref{fig:training_curve_global_nb2_h6}. As can be seen, the validation loss curve stabilizes early (around epoch 10) without diverging, indicating that the chosen hyperparameters, such as a dropout rate of 0.2 and a patience of 12, effectively control possible overfitting despite the high inter-subject variability inherent in the employed dataset.

\begin{figure}
\centering
\includegraphics[width=\columnwidth]{fig5_global_training_curves_nb2_h6.png}
\caption{Average training and validation loss across all five folds for the final CNN-Transformer configuration.}
\label{fig:training_curve_global_nb2_h6}
%\vspace{-10pt}
\end{figure}

% =============================
% TABLE: GRID SEARCH RESULTS
% =============================
\begin{table*}[t]
\centering
\small
\caption{Results of the statistical fine-tuning on the five-fold cross-validation.}
\label{tab:grid_search_results}
\footnotesize
\begin{threeparttable}
\resizebox{0.8\textwidth}{!}{
\begin{tabular}{l c c c c c c c}
\toprule
Config ID & LR & Drop & $L$ & Params & Mean F1 (SD) & Best Fold F1 & $p$-value$^{\dagger}$ \\
\midrule
Config-01 & 5e-4 & 0.1 & 1 & 232,290 & 0.8341 (0.026) & 0.8616 & 0.89 \\
Config-03 & 5e-4 & 0.2 & 1 & 232,290 & 0.8340 (0.022) & 0.8582 & 0.89 \\
Config-05 & 1e-3 & 0.1 & 1 & 232,290 & 0.8324 (0.023) & 0.8607 & 0.68 \\
\rowcolor[gray]{0.9} 
\textbf{Config-07} & \textbf{1e-3} & \textbf{0.2} & \textbf{1} & \textbf{232,290} & \textbf{0.8352 (0.015)} & \textbf{0.8524} & -- \\
Config-08 & 1e-3 & 0.2 & 2 & 399,762 & 0.8311 (0.022) & 0.8619 & 0.46 \\
Config-04 & 5e-4 & 0.2 & 2 & 399,762 & 0.8309 (0.019) & 0.8522 & 0.22 \\
Config-06 & 1e-3 & 0.1 & 2 & 399,762 & 0.8294 (0.020) & 0.8510 & 0.08 \\
Config-02 & 5e-4 & 0.1 & 2 & 399,762 & 0.8257 (0.021) & 0.8476 & 0.08 \\
\bottomrule
\end{tabular}
}
\begin{tablenotes}
\begin{minipage}{\textwidth}
    \item $^{\dagger}$ Two-sided Wilcoxon signed-rank test against \textit{Config-07} ($N=5$). $L$: Encoder layers. Fixed parameters: $d=144$, Heads=6, DW-Blocks=2, Warmup=4, Patience=12. The bold line denotes the best-selected configuration.
\end{minipage}
\end{tablenotes}
\end{threeparttable}
\vspace{-10pt}
\end{table*}

% -----------------------------------------------------------------------
% SUBSECTION: FINAL RESULTS (Updated Values)
% -----------------------------------------------------------------------
\subsection{Cross-subject test performance}

The winning model from the hyperparameter optimization phase, designated as \textit{Config-07}, which is built upon the \textit{nb2\_h6} backbone with a single Transformer encoder layer, a learning rate of $1\times 10^{-3}$, and a dropout rate of 0.2, demonstrates stable and competitive cross-subject performance. As evidenced by the fine-tuning results, this optimally configured CNN-Transformer achieves a mean macro F1-score of $0.835$, with a low standard deviation of $0.015$ across folds. This represents a solid improvement over the unoptimized baseline architecture, confirming the effectiveness of the hyperparameter optimization stage.

Related to the performance across folds, the best-selected model remained consistent, with macro F1-score values ranging from approximately $0.814$ in fold 3 to $0.852$ in folds 1 and 5. The worst-case scenario was observed in Fold 3 with a macro F1-score of $0.814$, highlighting the lowest-performing split, which aligns with prior analyses of the PhysioNet MI dataset, in which certain subject subsets are known to exhibit weaker or less consistent ERD/ERS patterns. Conversely, the model reaches peak performance in folds 1 and 5, reflecting clearer sensorimotor modulations among the subjects in those test splits.

Moreover, the overall confusion matrix across all five folds further characterizes the model's behavior as shown in Fig.~\ref{fig:confusion_global_allfolds_nb2_h6}. Misclassifications are remarkably balanced between L and R imagery, reaching 396 vs. 401 false predictions, respectively. This confirms that the model maintains a symmetric decision boundary without favoring either class, successfully delivering subject-independent decoding in an eight-channel setting while maintaining highly competitive efficiency without compromising classification performance. 

Overall, these results indicate that the proposed CNN-Transformer achieves stable subject-independent decoding under the specific inter-subject variability and recording conditions represented in PhysioNet. Although additional cross-dataset validation remains necessary to establish broader generalizability, the present findings provide strong evidence that compact hybrid architectures can remain competitive in low-density motor imagery decoding without calibration.

\begin{figure}
\centering
\includegraphics[width=0.6\columnwidth]{fig6_global_confusion_matrix_nb2_h6.png}
\caption{Overall confusion matrix accumulated across all five folds in the binary L/R task. Errors are balanced across both classes, indicating a symmetric decision boundary.}
\label{fig:confusion_global_allfolds_nb2_h6}
%\vspace{-10pt}
\end{figure}

\subsection{Ablation study}
\label{subsec:ablation}

To validate the contribution of specific architectural components to the overall performance, we conducted an ablation study on the proposed CNN-Transformer model. We compared the full baseline architecture against four variants: (i) removing the positional encodings to test temporal dependency modeling, (ii) replacing the learnable \texttt{[CLS]} token with global average pooling, (iii) substituting GroupNorm with standard BatchNorm, and (iv) using ReLU instead of ELU activations. All variants were trained using the same subject-wise five-fold cross-validation protocol. The obtained results derived from this process are summarized in Table~\ref{tab:ablation_study}.

\begin{table}[t]
\centering
\caption{Ablation study of architectural components across five folds.}
\label{tab:ablation_study}
\begin{threeparttable}
\footnotesize
\resizebox{\columnwidth}{!}{
\begin{tabular}{l c c c}
\toprule
Model Variant & Mean macro F1$\pm$SD  & $\Delta$ F1 & $p$-value$^{\dagger}$ \\
\midrule
(i) w/o Positional Encoding & 0.761 $\pm$ 0.024 & $-0.071$ & 0.06 \\
(ii) w/o \texttt{[CLS]} Token & 0.826 $\pm$ 0.019 & $-0.006$ & 0.31 \\
(iii) w/ BatchNorm & 0.822 $\pm$ 0.026 & $-0.010$ & 0.12 \\
(iv) w/ ReLU & 0.825 $\pm$ 0.018 & $-0.007$ & 0.31 \\
\rowcolor[gray]{0.9}
\textbf{Proposed method} & \textbf{0.835 $\pm$ 0.015} & -- & -- \\
\bottomrule
\end{tabular}
}
\begin{tablenotes}
\begin{minipage}{\columnwidth}
    \item $^{\dagger}$ Wilcoxon signed-rank test against the proposed model ($N=5$). $\Delta$ F1 denotes the absolute change in macro F1-score.
\end{minipage}
\end{tablenotes}
\end{threeparttable}
\vspace{-10pt}
\end{table}

From this table, it is evident that the most substantial impact was observed when removing positional encodings, resulting in a $7.4\%$ performance drop. Although this difference is marginally significant ($p=0.06$) due to the fold-wise sample size, the magnitude of the degradation strongly supports the conclusion that the Transformer encoder actively leverages the temporal order of EEG features to discriminate motor imagery patterns, rather than treating the signal as a bag-of-features.

Regarding feature aggregation, the use of a learnable \texttt{[CLS]} token consistently outperformed global average pooling ($\Delta \text{macro F1-score} = -0.009$). This suggests that the attention mechanism benefits from a dedicated token to dynamically aggregate global information, rather than relying on a static average over all time steps.

Finally, the choice of normalization and activation functions also contributed to stability. Replacing GroupNorm with BatchNorm degraded performance ($\Delta \text{macro F1-score} = -0.013$), likely due to the small batch sizes and high variance typical of subject-wise EEG training, where GroupNorm is known to be more reasonable. Similarly, ELU activations offered a slight advantage over ReLU ($\Delta \text{macro F1-score} = -0.010$), potentially by preserving negative signal fluctuations that contain relevant phase information in sensorimotor rhythms.

\subsection{Interpretability}
Understanding the decision mechanisms of the proposed CNN–Transformer model is important for assessing whether its predictions rely on physiologically meaningful EEG patterns. We conducted this analysis by considering four complementary perspectives: channel-level saliency, the statistical separability of EEG signals, the geometry of the learned feature space, and the temporal structure revealed by Transformer self-attention. These analyses allow us to verify whether the network relies on physiologically meaningful motor-imagery patterns, following established practices in EEG model interpretability \citep{tibor2017deep,nagarajan2023relevance,saibene2024deep}.

The global saliency map was quantified using the $Input \times Gradient$ method, defined as $|\nabla_x f(x) \odot x|$, where $f(x)$ is the model's output logit for the target class and $\odot$ denotes the element-wise (Hadamard) product. This measure was averaged across time and correctly classified test trials, and normalized to $[0,1]$. The resulting topographic saliency map across five folds is shown in Fig~\ref{fig:saliency_topomaps}. This map indicates that saliency is strongly concentrated in the bilateral sensorimotor regions (C3/FC3 and C4/FC4), whereas posterior channels contribute minimally. This distribution result aligns with the somatotopic organization of motor cortex and the contralateral $\mu/\beta$ desynchronization characteristic of hand imagery \citep{1niedermeyer2005electroencephalography,2pfurtscheller1999event,4pfurtscheller2001functional}, confirming that the model prioritizes relevant motor areas despite the reduced montage.

\begin{figure}
    \centering
    \includegraphics[width=0.7\linewidth]{fig7_global_saliency_map_nb2_h6.png}
    \caption{Global saliency map per channel (mean across five folds). The model focuses on the motor cortex (C3, C4), aligning with physiological expectations.}
    \label{fig:saliency_topomaps}
   % \vspace{-10pt}
\end{figure}

Moreover, to assess the intrinsic difficulty of the decoding task independently of the model, statistical tests were computed over the same eight-channel montage. A t-test between L and R imagery reveals a strongly lateralized pattern over C3 and C4, as shown in Fig.~\ref{fig:pvalues_comparison}. This result is consistent with distinct contralateral modulation. The high values of~$-\log_{10}(p)$ indicate reliable physiological differences, which help explain the successful classification performance attained by the proposed CNN-Transformer.

\begin{figure}
    \centering
    \includegraphics[width=0.7\linewidth]{fig8_global_log10p_map_nb2_h6.png}
    \caption{The $-\log_{10}(p)$ map quantifying intrinsic L/R EEG separability for the eight-channel montage.}
    \label{fig:pvalues_comparison}
   % \vspace{-10pt}
\end{figure}

Additionally, to visualize how the model organizes its learned representations, the \texttt{[CLS]} embeddings were projected onto a t-SNE plot (see Fig.~\ref{fig:tsne_binary}). From this plot, two clearly separated lobes emerge, corresponding to L and R imagery. This distinct clustering reflects the strong lateralized structure observed in the statistical maps and confirms that the Transformer encoder successfully maps the temporal EEG dynamics into a linearly separable latent space. Furthermore, the self-attention matrices from the final encoder layer of the proposed CNN-Transformer were averaged over heads, and the resulting values, which relate to how the Transformer integrates information across the 6s analysis window, are visualized in Fig.~\ref{fig:attention_binary}. As can be seen, the attention is concentrated in specific compact temporal regions rather than being uniformly distributed. This indicates that the model learns to identify and prioritize stable, discriminative time intervals, likely corresponding to the most prominent ERD/ERS events, while ignoring less relevant segments of the trial.

\begin{figure}
    \centering
    \includegraphics[width=\linewidth]{fig9_tsne_cls_fold3_nb2_h6.png}
    \caption{t-SNE projection of the \texttt{[CLS]} embeddings. The clear separation between L (blue) and R (orange) clusters indicates that the model learns discriminative features.}
    \label{fig:tsne_binary}
  %  \vspace{-10pt}
\end{figure}

\begin{figure}
    \centering
    \includegraphics[width=0.8\columnwidth]{fig10_global_attention_map_nb2_h6.png}
    \caption{Self-attention map averaged across heads. The vertical bands indicate that the model attends to specific discriminative time intervals within the 6-second window, which physiologically correspond to the task cue onset and peak ERD/ERS modulations.}
    \label{fig:attention_binary}
   % \vspace{-10pt}
\end{figure}

Across all interpretability analyses, the model consistently relies on neurophysiologically plausible cues. It focuses on sensorimotor channels, assigns relevance to time intervals consistent with motor imagery dynamics \citep{2pfurtscheller1999event,4pfurtscheller2001functional}, and learns feature representations in which lateralized classes remain highly structured. These findings provide a coherent interpretation of the quantitative results and confirm that the proposed lightweight architecture effectively captures the underlying neural mechanisms of motor imagery.

\subsection{Computational cost analysis}
The computational profile of the proposed CNN-Transformer highlights a potential solution for deployment on resource-constrained BCI hardware. The convolutional stem and Transformer encoder together require 0.024~GFLOPs per trial, with a total parameter count of 0.232~M. A detailed breakdown of these measurements is provided in Table~\ref{tab:comp_cost_cnn_tf}.

As detailed in the table, a full subject-independent fold (60 epochs, batch size 64) took 18.82~s to train, facilitating rapid experimentation and potential on-device recalibration. On the other hand, the test-time evaluation for a full fold took 8.69~s. Crucially, while the measured forward-pass latency of approximately 1.2 ms per batch of size B = 8 (i.e., approximately 0.15 ms per single trial) was obtained on a high-end GPU platform, the compact footprint of 0.232 M parameters and 0.024 GFLOPs suggests that the model is a promising candidate for future evaluation on clinical CPUs and embedded edge hardware. In this sense, the reported computational profile is consistent with the low-latency requirements ($<100$~ms) typically associated with online MI-BCI systems~\citep{saibene2024deep,wang2020accurate}, although direct benchmarking on limited-resource devices remains necessary. 

Regarding memory usage, the training process required 3281.2~MB of GPU memory, while inference used 159.5~MB. These values place the proposed architecture in an efficiency range comparable to compact EEG-specific models such as EEGNet \citep{9lawhern2018eegnet,wang2020accurate}, whereas providing the additional representational flexibility of a Transformer encoder. Overall, the proposed CNN-Transformer offers a practical accuracy–efficiency trade-off, remaining well within the limits of modern embedded or wearable BCI hardware \citep{george2022data,nagarajan2023relevance}.

\begin{table}[h]
\centering
\caption{Mean computational metrics of the proposed CNN-Transformer across five folds.}
\label{tab:comp_cost_cnn_tf}
\footnotesize
\begin{threeparttable}
\footnotesize
\resizebox{\columnwidth}{!}{
\begin{tabular}{lcccc}
\toprule
Phase & Time (s) & Memory (MB) & FLOPs (G) & Parameters (M) \\
\midrule
Model specs & --    & --     & 0.024 & 0.232 \\
Train       & 18.82 & 3281.2 & --    & --    \\
Test        &  8.69 & 159.5  & --    & --    \\
\bottomrule
\end{tabular}
}
\end{threeparttable}
\vspace{-10pt}
\end{table}

\subsection{Clinical impact and translational value}
\label{subsec:clinical_impact}

The proposed CNN–Transformer is designed for translational MI-BCI deployment. By achieving stable subject-independent decoding with a low-density eight-channel sensorimotor montage, the framework reduces dependence on dense EEG setups that typically increase preparation time and operational complexity in rehabilitation-oriented workflows. In addition, its compact computational profile (0.232~M parameters, 0.024~GFLOPs, and 159.5~MB test memory) supports the feasibility of future implementation on lower-resource hardware than that required by larger deep-learning pipelines. Although the present study was conducted offline on a high-end GPU platform and does not yet include embedded-device benchmarking, online closed-loop experiments, or patient validation, the reported results provide a technically plausible basis that may be relevant to future neurorehabilitation-oriented BCI systems. Although the results should not be interpreted as evidence of clinical efficacy, they demonstrated that the combination of low-density acquisition, calibration-free operation under the adopted protocol, and low algorithmic complexity can achieve promising subject-independent performance, making the proposed model an engineering basis for future evaluation in real-time wearable and bedside MI-BCI applications.

\subsection{Comparison with the state of the art}

\begin{table*}[t!]
\centering
\caption{Comparison with representative binary motor imagery decoding methods.}
\label{tab:sota_binary}
\footnotesize
\begin{threeparttable}
% Eliminamos el resizebox por completo.
% Usamos columnas 'p' para forzar el salto de línea en "Method" y "Notes".
\begin{tabular}{@{} >{\raggedright\arraybackslash}p{4cm} c c c c c c >{\raggedright\arraybackslash}p{3.5cm} @{}}
\toprule
Method & Dataset & Subjects & Montage & Channels & \makecell{Acc \\ (Global)} & \makecell{Acc \\ (FT / Adapt.)} & Notes \\
\midrule
EEG-Inception \citep{zhang2021eeg} & BCI IV-2b & 9 & Reduced & 3 & 77.5\% & 88.6\% & Subject-specific tuning \\
\addlinespace
Calibration-free TTA \citep{wimpff2024calibration} & BCI IV-2b & 9 & Reduced & 3 & 78.8\% & 83.5\% & Test-time adaptation \\
\addlinespace
Relevance CNN \citep{nagarajan2023relevance} & \multirow{2}{*}{KU MI} & \multirow{2}{*}{54} & Dense & 62 & 83.6\% & -- & Full montage \\
 & & & Reduced & 10 & 79.9\% & -- & Best 10 channels \\
\addlinespace
Transfer-learning CNN \citep{alimardani2023end} & Proprietary & 55 & Reduced & 16 & 62.5\% & -- & Calibration-free \\
\addlinespace
EEGNet \citep{lawhern2018eegnet} (Baseline) & PhysioNet & 103 & Reduced & 8 & 81.5\% & -- & Re-implemented under identical conditions \\
\addlinespace
CNN-only (Baseline) & PhysioNet & 103 & Reduced & 8 & 75.3\% & -- & Ablated proposed model \\
\addlinespace
\rowcolor[gray]{0.9}
\textbf{Proposed CNN-Transformer} & \textbf{PhysioNet} & \textbf{103} & \textbf{Reduced} & \textbf{8} & \textbf{83.3\%} & \textbf{--} & \textbf{Cross-subject} \\
\bottomrule
\end{tabular}
\begin{tablenotes}
    \item \textit{Note:} MI denotes motor imagery. Acc refers to classification accuracy. Reported accuracies correspond to cross-subject evaluation unless stated otherwise.
\end{tablenotes}
\end{threeparttable}
\end{table*}


A direct comparison with strictly controlled baseline architectures is provided in Table~\ref{tab:sota_binary} to evaluate the intrinsic capability of the proposed model. Under identical subject-wise cross-validation conditions, the proposed CNN-Transformer achieved an accuracy of 83.3\%, outperforming both the standard EEGNet \citep{lawhern2018eegnet} (81.5\%) and the CNN-only variant (75.3\%). The substantial performance drop observed when ablating the Transformer encoder (from 83.3\% to 75.3\%) demonstrates that relying solely on convolutional features is insufficient to capture the complex inter-subject variability present in low-density EEG. By integrating temporal self-attention, the proposed hybrid architecture successfully models long-range temporal dependencies, thereby achieving a competitive advantage and greater robustness across unseen subjects.

Furthermore, Table~\ref{tab:sota_binary} contextualizes the proposed model against other representative cross-subject decoding methods from the literature. It is important to note that direct comparisons across different studies are often complicated by variations in cohort size, montage density, and validation protocols. However, a general trend indicates that high performance in prior works is frequently associated with either dense spatial sampling or subject-specific adaptation. For instance, while Relevance CNN \citep{nagarajan2023relevance} reports an accuracy of 83.6\% using a high-density 62-channel montage, its performance drops to 79.9\% under a reduced 10-channel configuration. In contrast, methods designed for low-channel settings, such as EEG-Inception~\citep{zhang2021eeg} and calibration-free test-time adaptation approaches~\citep{wimpff2024calibration}, reach competitive accuracies (88.6\% and 83.5\%, respectively) but rely heavily on subject-specific fine-tuning or online adaptation on much smaller cohorts. Considering these factors, the proposed CNN-Transformer offers a highly favorable balance. Using only 8 EEG channels on a large 103-subject cohort, the model achieves 83.3\% global cross-subject accuracy without requiring any subject-specific calibration or fine-tuning, demonstrating strong potential for practical plug-and-play BCI deployment scenarios.

From a modeling perspective, this performance can be attributed to the complementary roles of the convolutional and Transformer components. The depthwise convolutional front-end provides parameter-efficient local feature extraction, whereas temporal self-attention aggregates long-range temporal dependencies and selectively weights informative patterns across time. This combination appears to mitigate, at least in part, the loss of spatial resolution inherent to sparse montages, hence preserving discriminative information even in low-density acquisition settings.

Finally, the proposed framework also offers clear advantages in computational efficiency. In contrast to calibration-dependent pipelines or larger deep architectures, the model remains compact, with only 0.232 million parameters and 0.024 GFLOPs per forward pass. Therefore, beyond its competitive decoding accuracy, the proposed CNN-Transformer constitutes a practical solution for low-density EEG-based MI decoding, combining cross-subject generalization, calibration-free operation, and low computational cost in a manner well aligned with real-world BCI deployment constraints.

\subsection{Study limitations}
Although the present study provides preliminary evidence under a strict subject-independent evaluation protocol, the following limitations should be noted:
\begin{itemize}  
    \item The proposed model was validated on a single public dataset, and additional cross-dataset evaluation will be necessary to assess robustness across acquisition systems, recording protocols, and broader sources of inter-subject variability.
    \item All experiments were conducted offline. Therefore, the implications for real-time closed-loop BCI operation remain to be validated experimentally.
    \item The present analysis is based on motor imagery recordings from healthy subjects and does not include patient cohorts, online rehabilitation scenarios, or embedded-device benchmarking. Accordingly, the translational implications of the proposed framework should be interpreted as engineering motivation for future work rather than as demonstrated clinical efficacy.
    \item The computational efficiency was estimated on a high-performance GPU. Thus, direct benchmarking on CPUs, embedded processors, and wearable platforms is still required. 
    \item  Although the proposed model was compared with representative literature, future work should include broader controlled comparisons under identical preprocessing and subject-wise splits.
\end{itemize}

%==============================================================
\section{Conclusions}
\label{sec:conclusions}
%==============================================================

This work demonstrates that a compact CNN–Transformer architecture can provide a favorable trade-off among low-density EEG decoding performance, subject-independent generalization, and computational efficiency in MI-BCIs. Under a strict subject-wise five-fold cross-validation protocol on the PhysioNet EEGMMI database, the optimized model achieved a mean global ACC of $0.833 \pm 0.014$ and a mean macro F1-score of $0.835 \pm 0.015$, demonstrating consistent subject-independent generalization despite the substantial inter-subject variability typical of EEG signals. These results are particularly significant given the use of a reduced eight-channel montage and a strict subject-independent evaluation protocol, both of which pose substantial challenges for MI decoding. An appropriate statistical analysis confirmed that larger architectural variants do not offer significant performance gains ($p > 0.05$), validating the design choice of a compact single-layer Transformer encoder. The proposed architecture integrates a convolutional front-end to extract localized spatiotemporal structures and self-attention mechanisms to capture long-range temporal dependencies, thus enabling efficient decoding of motor-related oscillatory dynamics. Furthermore, the model is computationally efficient, with approximately 0.232 M parameters and 0.024 GFLOPs per forward pass, supporting low-latency inference in the offline setting considered in this study and motivating future evaluation on embedded or wearable BCI platforms. This computational profile reinforces the practical relevance of the proposed architecture in scenarios where resource constraints and rapid inference are important design considerations.

Future work will extend this analysis toward real-world applicability by evaluating performance stability across wearable devices (e.g., BrainBit Flex~8) and exploring self-supervised pretraining, domain alignment, among other techniques, to reduce calibration demands in practical clinical settings.

\section*{Ethics statement}
This study used the publicly available EEG Motor Movement/Imagery dataset from PhysioNet. The data are de-identified and were collected independently of the present study. No new human-subject data were collected; therefore, additional institutional ethics approval was not required for this secondary analysis.

\section*{Data and code availability}
The raw EEG data are publicly available through the PhysioNet EEG Motor Movement/Imagery database. The derived subject splits, preprocessing scripts, model configuration files, and evaluation code are available at \url{https://github.com/ASP-ML/cnn_transformer_hybrid}. If required for peer review, an anonymized repository link will be provided.

\section*{Declaration of competing interest}
The authors declare that they have no known competing financial interests or personal relationships that could have appeared to influence the work reported in this paper.

\section*{Funding}
This work was supported by Universidad San Francisco de Quito USFQ through the Poli-Grants Program under Grants 33609 and 41993.

%\section*{Use of generative AI}
%During the preparation of this work, the authors did not use generative AI or AI-assisted tools to create, modify, or enhance scientific images. The authors reviewed and edited all manuscript content and take full responsibility for the final version.

\section*{Acknowledgments}
The authors thank the USFQ Applied Signal Processing and Machine Learning (ASP\&ML) Research Group for providing the computing infrastructure to implement and execute the developed source code. 

% To print the credit authorship contribution details
\printcredits
\balance
%% Loading bibliography style file
\bibliographystyle{cas-model2-names}

% Loading bibliography database
\bibliography{biblio}

\end{document}